\newcommand{\EsixText}{%
\textbf{Result.} Growing the masks by one token closes most of the gap: at $K{=}144$, the dilated oracle's adversarial leakage (\EsixLc) is statistically indistinguishable from the learned score's
(difference $-0.03$, 95\% CI $[-0.06,+0.01]$), and at $K{=}288$ it falls from 0.49 to \EsixLa, though it remains above the learned score (difference $+0.08$, CI $[+0.05,+0.12]$).
So a large part of the leakage really comes from tokens \emph{adjacent} to the annotated regions, which confirms the context explanation.
Growing by two tokens is \emph{worse} than by one (\EsixLb\ at $K{=}288$). The reason is the budget. Two-token dilation flags 300 of 576 tokens per image on average,
and in 56\% of images more than $K{=}288$, so selection must keep some flagged tokens. Dilation gives the ring the same score as the core, so the selector cannot prefer ring tokens over core tokens.
The learned score is \emph{graded} (mean 0.89 inside a mask, 0.47 in the first ring, 0.34 in the second), so under a tight budget it gives up the least risky tokens first.
\emph{Lesson for the second half:} privacy risk should be a graded, context-aware score, not a binary region. This favours the learned head and motivates distance-weighted oracles such as $\Pr_i=\max_j \Pr^{GT}_j\,\gamma^{d(i,j)}$ as a stronger reference.}
